% fig-d2-recognition-template.tex --- D2, the recognition template.
%
% Defines \ClebschRecognitionTemplate.
%
% The shape is drawn once and the table instantiates its two slots per
% paper.  Every pair below is taken from the named paper's abstract:
%   I    clebsch-rigidity/clebsch_rigidity.tex
%   II   clebsch-factorization/clebsch_factorization.tex
%   III  clebsch-passages/clebsch_passages.tex
%   IV   q13-passant-code/passant_code_q13.tex
%   AME  ame_lu (stabilizer AME local-unitary rigidity)

\providecommand{\ClebschRecognitionTemplate}{%
\begin{tikzpicture}[every node/.style={outer sep=0pt}]

\node[csfcentre,text width=36mm] (rich) at (-3.4,0)
  {rich object\\[1pt] the paper's full structure};
\node[csfnode,text width=36mm] (coarse) at (3.4,0)
  {coarse observation\\[1pt] low-degree or local data};

\draw[csfarr] (rich) to[bend left=22]
  node[csflab,above] {forget} (coarse);
\draw[csfarr,line width=1pt] (coarse) to[bend left=22]
  node[csflab,below] {recover (theorem)} (rich);

\node[anchor=north,inner sep=0pt] at (0,-1.9) {%
  \small
  \begin{tabular}{@{}l
    >{\raggedright\arraybackslash}p{40mm}
    >{\raggedright\arraybackslash}p{45mm}@{}}
  \toprule
  & \emph{rich object} & \emph{coarse observation}\\
  \midrule
  Paper I
    & the Clebsch code, a six-arc in \(\PGtwo{11}\)
    & its projective deep-hole syndrome locus\\\addlinespace[3pt]
  Paper II
    & the icosahedral conic matching and its two sheets
    & a two-valued strength-two quadratic trade\\\addlinespace[3pt]
  Paper III
    & the order-six conference operator \(C\)
    & four-local principal-minor data (the aligned four-sets)\\\addlinespace[3pt]
  Paper IV
    & the marked conic plane \(\PGtwo{13}\)
    & the minimum-word layer of its passant code\\\addlinespace[3pt]
  AME paper
    & a stabilizer \(\operatorname{AME}(2m,q)\) state
    & its product-unitary symmetries\\
  \bottomrule
  \end{tabular}};

\end{tikzpicture}%
}
